\section{H1: Activity in the settlement window}\label{sec:activity}

\subsection{Volume across the year}\label{sec:window}

Table~\ref{tab:activity} reports how the settlement window on each type of session differs
from the same security's window on ordinary sessions. On an ordinary session the window
accounts for about \ActWindowshareLiquidNormal{}\% of the session's volume in the liquid
group, and similar shares in the other two.

\begin{table}[htbp]
\centering
\small
\caption{Activity in the settlement window by type of session. Each entry is the coefficient
on an indicator for the session type in a regression with security fixed effects, estimated
separately for each group, and measures the difference from ordinary sessions of the same
security. The last column is the difference between the securities with derivatives and the
placebo group, estimated with security and session fixed effects. Standard errors clustered
by session in parentheses. $^{*}$, $^{**}$ and $^{***}$ denote significance at the ten, five
and one per cent levels.}
\label{tab:activity}
\resizebox{\textwidth}{!}{\input{generated/tables/activity}}
\end{table}

On monthly expiries, trading is drawn into the window in securities with derivatives and not
in the others. The window's share of session volume rises by \ActWindowshareLiquidMonthlyexpiry{}
percentage points in the liquid group and by \ActWindowshareIlliquidMonthlyexpiry{} points
in the illiquid group, while the change in the placebo group, \ActWindowsharePlaceboMonthlyexpiry{}
points, is not distinguishable from zero. The difference between the derivatives groups and
the placebo group is \ActWindowshareDidMonthlyexpiry{} points, with a standard error of
\ActWindowshareDidMonthlyexpirySe{}. The calendar controls do not reproduce this pattern. On weekly
index expiries the window's share of volume does not change in any group. On month ends it
rises sharply in all three groups, by \ActWindowshareLiquidMonthend{} points in the liquid
group and by \ActWindowsharePlaceboMonthend{} points in the placebo group, which reflects the
closing demand of index and portfolio rebalancing; the additional rise in the securities with
derivatives, \ActWindowshareDidMonthend{} points, is smaller than on monthly expiries.

The settlement price also moves further from the pre-window price on monthly expiries. For
the securities with derivatives the absolute move increases by
\ActAbsdriftDerivativeMonthlyexpiry{} basis points from an average of
\ActAbsdriftDerivativeNormal{}, and the final five minutes of the window end
\ActAbsterminalDerivativeMonthlyexpiry{} basis points further from the settlement price.
Neither increase is specific to the derivatives groups. The placebo group shows increases of
similar size, and the differences between the two, \ActAbsdriftDidMonthlyexpiry{} and
\ActAbsterminalDidMonthlyexpiry{} basis points, are small and insignificant. Figure~\ref{fig:daytypes}
shows the same comparison in levels.

\begin{figure}[htbp]
  \centering
  \includegraphics[width=\textwidth]{day_types.pdf}
  \caption{Share of session volume traded in the settlement window, and the absolute move
    from the pre-window price to the settlement price, by group and type of session. Means
    over security-sessions.}
  \label{fig:daytypes}
\end{figure}

\subsection{The order book}\label{sec:book}

On the \DesignSessions{} sessions of the order book sample the limit order book is
reconstructed, which shows how the supply of liquidity in the window changes on expiry days.
Each measure is computed per security and session and compared between each monthly expiry
and the control session in the same month, the trading day before it.
Table~\ref{tab:descriptives} reports the averages and Table~\ref{tab:contrasts} compares the
expiry-minus-control differences of the securities with derivatives with those of the
placebo group. The full set of paired tests is reported in Appendix~\ref{app:tests}.

\input{generated/tables/descriptives}

The window is more expensive to trade in on expiry days. The quoted spread widens from
\SpreadMeanBpsControl{} to \SpreadMeanBpsExpiry{} basis points, the displayed notional in the
book falls from \BookVisibleCrControl{} to \BookVisibleCrExpiry{} crore, and sweeping the
displayed book with an order of \ProbeSizeCr{} crore moves the price by
\MoveBpsSmallExpiry{} basis points against \MoveBpsSmallControl{} on control sessions. The share
of orders cancelled within one second of entry rises from \LifePhantomRateControl{}\% to
\LifePhantomRateExpiry{}\%. The widening of the spread and the rise in fleeting orders are
specific to securities with derivatives: relative to the placebo group, the spread widens by an
additional \DidMeanspreadbpsPlacebo{} basis points ($p = \DidMeanspreadbpsPlaceboP{}$) and the
share of orders cancelled within a second rises by an additional
\DidPhantomorderratePlacebo{} ($p = \DidPhantomorderratePlaceboP{}$). The
cost of sweeping the book rises by a similar amount in both sets of securities.

The rise in cost is modest relative to the level. Sweeping the book is several times more
expensive in the illiquid group, at \MoveBpsSmallIlliquid{} basis points, than in the liquid
group, at \MoveBpsSmallLiquid{}, on any session. The cost of moving the price is set far more by
the security than by the day.

The share of resting size that is not displayed is high on every session, at
\ConcealedDepthShareControl{}\% on control sessions, and does not change on expiry. Concealed
orders become somewhat larger relative to displayed ones, from \ConcealedSizeMultipleControl{}
to \ConcealedSizeMultipleExpiry{} times their size, and are replenished slightly less often per share
traded. Fewer, larger concealed orders worked more slowly through the window describe the
execution of large expiring positions rather than an attempt to hide an accumulation aimed at
the settlement price.

\input{generated/tables/contrasts}
